% generated by roam.py handbook export -- не редактировать, править заметки
% 20261010095611-порядок_метода_оптимизации.org (body 0c279e92f336)

\begin{definition}[{Порядок метода оптимизации}]\label[definition]{note:950F2C19-C2C7-42CD-9CDF-0B9AC03A3DC8}
  Метод оптимизации называют методом \emph{$k$-го порядка}, если он использует производные $f$ до $k$-го порядка включительно:

  \begin{itemize}
    \item метод \emph{нулевого порядка}~--- только значения функции (например, \hyperref[note:96DED7D9-3918-45C5-9850-E05030941645]{метод золотого сечения}, \hyperref[note:4E950058-62DB-4EDA-A71D-98CACBBC41AA]{метод парабол});
    \item метод \emph{первого порядка}~--- значения и градиент (например, \hyperref[note:124BD9BC-084E-43FB-B617-AF8690B33EF8]{дихотомия по производной}, \hyperref[note:BEEE81AD-13FC-4846-A160-C72E31760BEE]{градиентный спуск});
    \item метод \emph{$k$-го порядка}~--- $f^{(k)}$.
  \end{itemize}
\end{definition}
